\scanpage{24}

d'où
\[
\partial_n:\underline H_Y^1(G_{mX})\longrightarrow\underline H_Y^2(\mu_{nX})
\]
et nous poserons
\[
\deleted{\unsure{\varphi}}\ 
\underline{\unsure{\gamma|\sigma}}_{Y/X}^{(n)}
=\partial_n(\unsure{\beta_{Y/X}})
\in\Gamma(X,\underline H_Y^2(\mu_{nX}))
\]
en laissant tomber les exposants $n$
s'il n'y a pas d'ambiguité possible.

\inserted{La définition de}
Cette section \deleted{satisfait} commute à la localisation
de Zariski (et $=$ à $\forall$ changement de base
$X'\longrightarrow X$ tel que $Y'$ reste une ``int. complète''
i.e. $\varphi'$ reste section régulière) ce qui
permet de définir $\unsure{\gamma|\sigma}_{Y/X}^{(n)}$ en général
(sans équation globale). Si \inserted{maintenant} $Y$ est intersection
complète de codim $d$, alors localt $Y$ est
de la forme
\[
Y=Y_1\cap\ldots\cap Y_d
\]
et on posera
\[
\underline{\unsure{\gamma|\sigma}}_{Y/X}
=\underline{\unsure{\gamma|\sigma}}_{Y_1/X}^{\deleted{\unsure{(n)}}}
\cdot\ldots\cdot\underline{\unsure{\gamma|\sigma}}_{Y_d/X}
\in\Gamma(Y,\underline H_Y^{2d}(\mu_n^{\otimes d}))
\]
(\uline{cup-produit}), et on prouve que cette
section ne dépend pas du choix des $Y_i$,
donc se définit globalement.

Propriétés générales :

\mspage{23}
